\begin{abstract}
A physical agent can get better over a few trials without learning a reusable model of the world. It may search locally, remember a useful trajectory, or exploit visual regularities in the tested scene. We study a stricter question: after budgeted interaction with one latent physical world, can a multimodal agent predict counterfactual future frames for new states governed by the same hidden dynamics? We introduce a benchmark built around this question. Each episode is indexed by a five-level hierarchy: physical-phenomenon family, world-parameter regime, latent-dynamics condition, latent-law realization, and counterfactual prediction instance. During support, the agent queries a tool-mediated interface over several instances. The interface returns observable state but withholds the latent law, hidden constants, and answer keys. During evaluation, tool access is frozen. The agent receives a new initial scene and four rendered candidate futures, whose positions are independently permuted, and must select the physically correct frame. The benchmark contains 5,670 instances across seven families and reports performance by family, regime, latent-dynamics condition, realization, and interaction budget. Baseline results on frontier multimodal agents, evaluated as complete model--harness systems, quantify how far current systems are from transferring physical evidence rather than fitting the interaction history. \TODO{Add the one-sentence quantitative headline once baseline runs finish.}
\end{abstract}
